{\huge\chapter{Заключение}}
\label{sec:Chapter4} \index{Chapter4}

В работе рассмотрена задача реидентификации транспортных средств по внешнему виду при съёмке с разноракурсных камер наблюдения. Построен воспроизводимый пайплайн, включающий предобработку, обучение эмбеддинг-моделей (ResNet-50 и ViT-Small с BNNeck, комбинированные функции потерь CE + Triplet) и этап сопоставления с постобработкой (k-reciprocal re-ranking). Эксперименты на наборе VeRi-776 демонстрируют, что корректная компоновка современных практик позволяет достигать высоких показателей качества без использования распознавания номерных знаков.

\section*{Основные результаты}
\begin{enumerate}
  \item Сформулированы теоретические основы Vehicle ReID и требования к инвариантности признаков при доменном сдвиге между камерами.
  \item Реализованы и сравнены два подхода: CNN (ResNet-50) и Transformer (ViT-Small) с BNNeck и комбинированной функцией потерь.
  \item ViT-Small показал лучшие результаты: Rank-1 90.0\%, mAP 69.4\% против Rank-1 88.7\%, mAP 64.3\% у ResNet-50 (разница +5.1\% по mAP, +1.3\% по Rank-1).
  \item Показано, что постобработка re-ranking стабильно повышает mAP на 5--7\% при умеренных издержках по времени инференса.
  \item Проведено абляционное исследование влияния аугментаций, триплетного обучения и BNNeck на итоговые метрики.
\end{enumerate}

\section*{Ограничения работы}
\begin{itemize}
  \item Использовались предоставленные кропы объектов; влияние качества детекции и трекинга не анализировалось.
  \item Эксперименты проведены на статичных изображениях; видеоинформация (темпоральный контекст) не использовалась.
  \item Тестирование выполнено на одном датасете VeRi-776; оценка на других бенчмарках остаётся направлением для расширения.
\end{itemize}

\section*{Направления дальнейших исследований}
\begin{itemize}
  \item Интеграция видео-контекста: темпоральные модели и трек-ассоциированный ReID (TCL, 3D-CNN/ViT).
  \item Domain Adaptation/Generalization между городами и оптикой (UDA/SSL, псевдометки, CamStyle/MixStyle, IBN).
  \item Эффективная индексация больших галерей (IVF-PQ/HNSW в FAISS) и сжатие эмбеддингов (OPQ/Distillation).
  \item Расширение объектов идентификации: велосипеды, автобусы и др.; мультидоменная оценка (VERI-Wild, CityFlowV2).
  \item Исследование более крупных Vision Transformer моделей (ViT-Base, ViT-Large) с предобучением на больших датасетах.
\end{itemize}

В целом, полученные результаты подтверждают применимость и устойчивость эмбеддинг-подходов реидентификации для практических задач городского наблюдения без опоры на распознавание номерных знаков. Сравнение CNN и Transformer архитектур показало значительное преимущество ViT-Small по качеству (+5.1\% mAP) благодаря механизму глобального внимания, при умеренном компромиссе в скорости инференса. Предложенный пайплайн может служить основой для дальнейших разработок и внедрения.
